Recent years have witnessed the deployment of high performance networks, e.g., FTTH and 5G mobile networks, to support the stringent and requirements in terms of latency, bandwidth, reliability, and jitter of emerging applications. For instance, remote surgery requires reliability and low-latency for video streaming and control feedback; AR/VR applications and the Metaverse require high-bandwidth and low-latency to allow human interaction with the environment. Cloud Gaming (CG) applications also face challenges in terms of delay and bandwidth consumption which are hardly met by current cellular network architectures, exposing CG users to network impairments that deteriorate their \ac{QoE}. \ac{ISP}s need efficient monitoring techniques to detect QoE deterioration that their customers may experience.

A primary approach widely adopted in networking to detect users' QoE degradation is to rely on \ac{AD} methods which are either applied manually or rely on rule-based techniques \cite{cannady1998artificial}. 
However, the increasing complexity of network infrastructures tend to make these techniques impractical and inefficient. The recent advances in Machine Learning (ML) have been leveraged in various domains, including networking where large amounts of data are available and could be used to build classification and prediction models. The popularity of ML is mainly due to the success of supervised learning which requires a plethora of labeled data during the training phase. However, data labeling has to be done by domain experts and has proven to be a long and tedious process. To circumvent the need for labeled data, unsupervised learning techniques are increasingly adopted, in particular for anomaly detection.

Many anomaly detection models based on ML techniques have been proposed in the literature, including distance-based algorithms \cite{liu2008isolation,Ramaswamy2000EfficientAF}, predictive algorithms, reconstruction-based algorithms \cite{zong2018deep, audibert2020usad}, one-class algorithms \cite{scholkopf1999support, pmlr-v80-ruff18a}, generative algorithms \cite{GoodfellowIan2020GenerativeAN}, etc. The performance of these models is usually evaluated using the F1-score as the key performance metric. The majority of these studies also used a \textit{point-wise} approach to compute the F1-score, i.e., they consider only anomalous observations. The issue with the \ac{PW} approach is that it disregards the fact that degradations can occur in the form of consecutive anomalous observations. To address this limitation, new approaches \cite{xu2018unsupervised,Lavin2015EvaluatingRA, hundman2018detecting} are proposed to better evaluate the performance of the AD models. These approaches consider degradations as \textit{windows} of anomalies, aggregate the predictions following different criteria and compute the F1-score accordingly. Another issue presented by many of existing ML models is the difficulty in comparing them for a given application since each model is evaluated using a different methodology, benchmark datasets and evaluation metrics. Each method reports higher F1-score than its competitors making it tricky for domain experts to choose the most appropriate approach that fits their needs.

In this paper, we present a comprehensive comparison and analysis on the performance of unsupervised ML models through experiments while relying on a consistent evaluation methodology. In particular, we focus on the detection of QoE degradation in low-latency applications using a real-world multivariate time-series dataset of Key Performance Indicators in CG sessions collected under different 4G network conditions. The game sessions are recorded using the public Google Stadia CG platform. We objectively evaluate the ability of each model to detect anomalies that can lead to QoE degradation for CG users. The experiments conducted under 4G conditions and the conclusions drawn from them remain applicable to 5G networks, as the QoE degradation we aim to detect are related to features collected at the CG application-level.
In addition, we provide insights and offer recommendations to network operators on the most appropriate AD model that meets their application requirements for anomaly detection, including real-time or offline inference, detection accuracy, robustness to data contamination rate due to different environments, etc. 

\contributions{
The main contributions of this chapter can be summarized as follows:
\begin{itemize}
    \item We demonstrate, based on synthetic models, the limitations of existing \textit{window} approaches for evaluating the performance of AD models and then propose our \ac{WAD} approach.
    \item We perform an exhaustive evaluation of ML models to assess (i) their robustness by injecting anomalies in their training sets (i.e., data contamination), and (ii) their capability to detect short and longer users' QoE degradation by varying the windows size.
    \item Based on our experiments, we offer recommendations to network operators and network management experts on the best models for different application requirements.

\end{itemize}
}


% The remainder of this paper is organized as follows. Section \ref{chap_3:related_works} presents the related work. Section \ref{chap_3:methodology} describes our methodology for anomaly detection and exposes the proposed WAD approach. Section \ref{chap_3:evaluation} compares the performance of the ML models and Section \ref{chap_3:discussion} discusses which model is best depending on the use case scenario. 